\section{Boundary averaging and measurable correction}\label{sec:boundary}

The finite-factor argument begins with nearby unitaries whose products
are only approximately respected.  For their countable indexing group
$G$, we construct a probability space $(D,\nu)$ with a nonsingular
$G$-action and an equivariant mean that corrects these multiplication
errors.  On $D^2$ two actions have different roles.  The separate action
$(g,h)(x,y)=(gx,hy)$ permits averaging of cocycles.  The diagonal action
$g(x,y)=(gx,gy)$ has a rigidity property: bounded measurable vector
fields transforming under a unitary representation are constant.
For cocycles with values in a finite algebra, a tracial convex-hull
construction will connect these two uses of $D^2$.  All estimates are
independent of $G$.

For a nonsingular action of a countable group $G$ on a probability space
$(D,\nu)$, write
\[
 \alpha_g F(x)=F(g^{-1}x),\qquad
 (\lambda_g F)(t,x)=F(g^{-1}t,g^{-1}x).
\]
Nonsingular means that the action preserves the class of null sets;
measure preservation is not required.  On $G\times D$ we use counting
measure in the first variable.  A bounded operator field on a separable
Hilbert space $K$ means an essentially bounded, weakly measurable map into
$\B(K)$.  Its norm is the essential supremum of its operator norms.
Separability makes this a measurable function: the operator norm can be
tested on countably many unit vectors.  Products of such fields are again
weakly measurable, since their values on fixed vectors are measurable
$K$-valued maps and can be approximated by countably valued maps.

The scalar mean below is the conditional-expectation formulation of
amenability for a measured group action
\cite{Zimmer1978Boundary,AdamsElliottGiordano1994}; see in particular
\cite[Theorem~3.4, p.~810]{AdamsElliottGiordano1994}.
We give the construction and the operator-field extension required here.

\begin{lemma}[Equivariant boundary means]\label{lem:boundary-mean}
For every countable group $G$ there are a standard probability space
$(D,\nu)$ with a nonsingular $G$-action and a positive unital contraction
\[
 m:L^\infty(G\times D)\longrightarrow L^\infty(D)
\]
satisfying
\begin{equation}\label{eq:boundary-mean-identities}
 m(\lambda_g F)=\alpha_g(mF),\qquad m(fF)=f\,m(F)
 \quad(f\in L^\infty(D)).
\end{equation}
There is a mean $m^{(2)}$ with the same properties for the separate
$G\times G$-action on $D\times D$ and functions on
$G^2\times D^2$.  Each mean extends to a contraction on bounded operator
fields on any separable $K$.  The extension preserves adjoints and
von Neumann algebra values, and is a bimodule map over bounded operator
fields independent of the group variable:
\begin{equation}\label{eq:boundary-operator-module}
 m(BFC)=B\,m(F)\,C.
\end{equation}
\end{lemma}

\begin{proof}
Choose a symmetric probability measure $\mu$ on $G$ with $\mu(g)>0$ for
every $g$.  Let $X=G^{\N\cup\{0\}}$ with product probability
$\mathbb P=\mu^{\N\cup\{0\}}$.  On $X$ define
\[
 g(\omega_0,\omega_1,\ldots)=(g\omega_0,\omega_1,\ldots),\qquad
 S(\omega_0,\omega_1,\ldots)=(\omega_0\omega_1,\omega_2,\ldots).
\]
These maps commute.  The first-coordinate laws of their pushforwards are
respectively $g\mu$ and $\mu*\mu$, while all later coordinates remain
independent with law $\mu$.  Full support therefore makes both
pushforwards equivalent to $\mathbb P$.

Let $\mathcal I$ be the completed sigma algebra of sets $E$ satisfying
$S^{-1}E=E$ modulo null sets.  It is countably generated modulo null sets.
Indeed the measure algebra of $X$ is separable for the metric
$\mathbb P(E\mathbin\triangle F)$, so its subspace $\mathcal I$ is
separable.  Choose a countable dense family $E_j$ in that subspace.
For each $E\in\mathcal I$, a sequence from this family with summable
symmetric-difference errors converges almost everywhere in indicators
to $1_E$.  Thus the $E_j$ generate $\mathcal I$ after completion.
The $G$-action preserves $\mathcal I$, because it commutes with $S$.
Include all translates $gE_j$ and set
\[
 b(\omega)_{j,g}=1_{gE_j}(\omega),\qquad
 D=\{0,1\}^{\N\times G},\qquad \nu=b_*\mathbb P.
\]
Equip $D$ with its Borel probability measure and, when taking measurable
functions, its completion.  The coordinate-permutation action
$(h x)_{j,g}=x_{j,h^{-1}g}$ satisfies $b(h\omega)=h b(\omega)$.
It is nonsingular, since the action on $X$ is nonsingular.  We have
$b\circ S=b$ almost everywhere.  Pullback by $b$ identifies
$L^\infty(D)$ with the $S$-fixed subalgebra of $L^\infty(X)$: one
inclusion follows from $b\circ S=b$, and the other from the generation
of $\mathcal I$.  This description also applies to scalar $L^2$
functions.  Countably many representative identities can always be
made simultaneous by deleting a null set.

For $F\in L^\infty(G\times D)$ put
\[
 (JF)(\omega)=F(\omega_0,b(\omega)),\qquad
 T_nF=\frac1n\sum_{j=0}^{n-1}(JF)\circ S^j.
\]
The map $J$ is well defined on equivalence classes: there are only
countably many group coordinates, and each section has a $\nu$-null
exceptional set.  These maps are positive unital contractions.  Take a
simultaneous pointwise weak-star cluster subnet of the maps $T_n$, using
the product of the weak-star compact balls
$\prod_F\{a\in L^\infty(X):\norm{a}\le\norm{F}\}$.  Denote its
limit by $T$; it is linear, positive, unital and contractive.

Pullback by a nonsingular measurable map is weak-star continuous on
$L^\infty$.  To see this directly, for $q\in L^1$ the measure
$E\mapsto\int_{S^{-1}E}q\,d\mathbb P$ is absolutely continuous with
respect to $\mathbb P$; its Radon--Nikodym density is the required
preadjoint applied to $q$.  Its $L^1$ norm is at most $\norm{q}_1$.
Consequently the telescoping bound
\[
 \norm{(T_nF)\circ S-T_nF}\le 2\norm{F}/n
\]
implies $(TF)\circ S=TF$.  Write $TF=(mF)\circ b$.
Moreover $J\lambda_g F=(JF)\circ g^{-1}$ and the group action commutes
with $S$.  Weak-star continuity of its nonsingular pullback gives the
first identity in \eqref{eq:boundary-mean-identities}.  The identity
$b\circ S=b$ gives $T_n(fF)=(f\circ b)T_nF$; multiplication by a fixed
bounded function is weak-star continuous, proving the second identity.

For the product mean use $X^2$, its two commuting shifts $S_1,S_2$,
and rectangular averages
\[
 \frac1{nm}\sum_{i=0}^{n-1}\sum_{j=0}^{m-1}
 F\bigl((S^i\omega)_0,(S^j\omega')_0,b(\omega),b(\omega')\bigr),
 \qquad n,m\longrightarrow\infty.
\]
A simultaneous weak-star cluster is fixed by both shifts.  These common
fixed functions are exactly the functions of $(b(\omega),b(\omega'))$.
For completeness, a fixed scalar $L^2(X^2)$ function has almost every
section in the closed subspace $L^2(\mathcal I)\subset L^2(X)$ in
each variable, by Fubini and the preceding identification.  The two
orthogonal projections onto
$L^2(\mathcal I)\otimes L^2(X)$ and
$L^2(X)\otimes L^2(\mathcal I)$ commute; their product has range
$L^2(\mathcal I)\otimes L^2(\mathcal I)$.  This proves the assertion.
The preceding equivariance and module arguments now apply in both
coordinates.

Here are details of the operator-field extension; they also explain why
no normality of $m$ is needed.  On a countable rational complex vector
subspace dense in $K$, define the coefficients of $m(F)(x)$ by applying
the scalar $m$ to $\ip{F(t,x)\xi}{\eta}$.  Outside a single null set
these coefficients are sesquilinear and bounded by
$\norm{F}\norm{\xi}\norm{\eta}$.  They therefore define a bounded
operator, and extension by density supplies all coefficients.
Positivity of the scalar mean implies that this extension preserves
adjoints; the norm bound gives contractivity.  The coefficient
identities also prove equivariance.

The scalar module identity implies the same coefficient formula when
the test vectors depend measurably on $x$.  First take countably valued
bounded vector fields: on each member of the countable joint partition,
localize using the scalar module identity and use the fixed-vector
formula.  Then approximate arbitrary bounded measurable vector fields
uniformly by countably valued ones, using separability of $K$ and the
norm bound.  Testing $BFC$ on fixed vectors with the variable vectors
$C(x)\xi$ and $B(x)^*\eta$ proves
\eqref{eq:boundary-operator-module}.  In particular, if $F$ takes values
in a von Neumann algebra $A\subset\B(K)$, its mean commutes with each
constant operator in $A'$.  Choose a countable strongly dense subset of
the unit ball of $A'$ to make this simultaneous almost everywhere;
then $m(F)(x)\in A''=A$.  Every step applies also to $m^{(2)}$.
\end{proof}

The next proof uses the averaging and polar-decomposition iteration
that appears in stability of approximate unitary representations
\cite{Kazhdan1982,BurgerOzawaThom2013}; compare
\cite[Theorem~3.2 and its proof, pp.~116--117]{BurgerOzawaThom2013}.
Here the mean belongs to an amenable action, and the corrected objects
are measurable cocycles.

\begin{lemma}[Correction of a measurable unitary cocycle]
\label{lem:boundary-correction}
Let $D$ and $m$ be as in \Cref{lem:boundary-mean}, and let
$A\subset\B(K)$ be a von Neumann algebra on a separable Hilbert space.
Suppose $V_g\in L^\infty(D,\U(A))$, $g\in G$, and
\[
 e=\sup_{g,h\in G}\norm{V_g\alpha_g(V_h)-V_{gh}}_\infty
 \le\frac1{64}.
\]
Then there are measurable unitary fields $W_g$ with
\begin{equation}\label{eq:boundary-exact-cocycle}
 W_g\alpha_g(W_h)=W_{gh},\qquad
 \sup_g\norm{W_g-V_g}_\infty\le 2e.
\end{equation}
In particular $W_1=I$.  The identities hold simultaneously almost
everywhere for every $g,h$, and may be used on a $G$-invariant conull
subset of $D$.
\end{lemma}

\begin{proof}
We give a quadratic improvement and then iterate it.  For the current
family define
\[
 H_g(t,x)=V_t(x)V_{g^{-1}t}(g^{-1}x)^*,\qquad A_g=m_t H_g.
\]
The defect inequality implies $\norm{H_g-V_g}_\infty\le e$, hence
$\norm{A_g-V_g}_\infty\le e$ and
$\norm{H_g-A_g}_\infty\le2e$.  All $H_g$ are unitary and all $A_g$
are contractions.  There is the exact factorization
\[
 H_{gh}=H_g\lambda_g(H_h),\qquad
 m(\lambda_g H_h)=\alpha_g(A_h).
\]
Using the left and right module identities to cancel the linear terms
in the product gives
\begin{align*}
 A_{gh}-A_g\alpha_g(A_h)
 &=m\bigl((H_g-A_g)(\lambda_g H_h-\alpha_g(A_h))\bigr),\\
 I-A_g^*A_g&=m\bigl((H_g-A_g)^*(H_g-A_g)\bigr).
\end{align*}
Both right sides have norm at most $4e^2$.  The same calculation gives
$\norm{I-A_gA_g^*}\le4e^2$.  Since $A_g$ is within $e<1$ of a
unitary, it is invertible in $L^\infty(D,A)$.  Its unitary polar part
$V'_g=A_g(A_g^*A_g)^{-1/2}$ satisfies
\[
 \norm{V'_g-A_g}_\infty
 =\norm{I-|A_g|}_\infty\le4e^2.
\]
Functional calculus preserves measurability: here it can be obtained
by uniform polynomial approximation on a fixed compact interval away
from zero.  The new defect is at most $16e^2$, since changing the three
terms in $A_g\alpha_g(A_h)-A_{gh}$ to their polar parts costs at most
$12e^2$.  Also
$\sup_g\norm{V'_g-V_g}_\infty\le e+4e^2$.

Starting with $e_0=e$, the successive defects satisfy
$e_{j+1}\le16e_j^2\le e_j/4$.  Thus the families converge uniformly in
$g$ and in essential operator norm to unitary fields $W_g$.  Their
defects tend to zero, and the total change is bounded by
\[
 \sum_{j\ge0}(e_j+4e_j^2)
 \le\frac{17}{16}\sum_{j\ge0}e_j
 \le\frac{17}{12}e\le2e.
\]
If $e=0$, the original family already works.  Countability of $G$
and nonsingularity permit deletion of all exceptional sets and their
translates.  Finally the cocycle identity with $g=h=1$ makes the
unitary $W_1$ idempotent, so $W_1=I$.
\end{proof}

\begin{lemma}[Averaged intertwiner]\label{lem:boundary-intertwiner}
Take either $(\Gamma,Z)=(G,D)$ or
$(\Gamma,Z)=(G^2,D^2)$, with its mean from
\Cref{lem:boundary-mean}.  Let $U:\Gamma\to\U(K)$ be a unitary
representation on a separable Hilbert space, and let
$\Sigma_s\in L^\infty(Z,\U(K))$ be an exact cocycle for the given
action.  If
\[
 \sup_{s\in\Gamma}\norm{\Sigma_s-U_s}_\infty\le\delta<1,
\]
there is $L\in L^\infty(Z,\U(K))$ such that
\begin{equation}\label{eq:boundary-intertwiner}
 \norm{L-I}_\infty\le2\delta,\qquad
 L(z)U_s=\Sigma_s(z)L(s^{-1}z).
\end{equation}
\end{lemma}

\begin{proof}
Set $B=m_t(\Sigma_tU_t^*)$.  Then $\norm{B-I}_\infty\le\delta$,
so $B$ is invertible.  Equivariance, the module identities and the
cocycle identity give
\[
 \Sigma_s\alpha_s(B)U_s^*
 =m_t\bigl(\Sigma_s\alpha_s(\Sigma_{s^{-1}t})
               U_{s^{-1}t}^*U_s^*\bigr)
 =m_t(\Sigma_tU_t^*)=B.
\]
Thus $BU_s=\Sigma_s\alpha_s(B)$.  Taking adjoints and products shows
$B^*B=U_s\alpha_s(B^*B)U_s^*$, so the same identity holds after
replacing $B$ by its unitary polar part $L$.  The inequalities
$(1-\delta)I\le|B|\le(1+\delta)I$ give
$\norm{L-B}_\infty\le\delta$, proving the stated bound.
\end{proof}

Applied on $D^2$, the preceding lemma gives an intertwiner for the
separate action.
For the finite-factor application we also need an algebra-valued
comparison of the two boundary variables under the diagonal action.
The next construction uses the trace in place of an invariant mean
for that action.

\begin{lemma}[Diagonal comparison in a finite algebra]
\label{lem:boundary-diagonal-unitary}
Let $A$ be a finite von Neumann algebra with faithful normal tracial
state $\tau$ and separable trace Hilbert space $L^2(A,\tau)$.
Let $W_g\in L^\infty(D,\U(A))$ be an exact cocycle.  Suppose
\[
 \sup_{g\in G}\mathop{\rm ess\,sup}_{x,y\in D}
 \norm{W_g(x)W_g(y)^*-I}\le r<1.
\]
Then there is a measurable field $a:D^2\to\U(A)$ such that
\begin{equation}\label{eq:boundary-diagonal-unitary}
 \norm{a-I}_\infty\le2r,\qquad
 a(x,y)=W_g(x)a(g^{-1}x,g^{-1}y)W_g(y)^*.
\end{equation}
In particular, if $\sup_g\norm{W_g-u_g}_\infty\le\delta$ for any
constant unitaries $u_g$ on a common representation space, the conclusion
holds with $\norm{a-I}_\infty\le4\delta$ when $2\delta<1$.
\end{lemma}

\begin{proof}
For almost every $(x,y)$ let $C_{x,y}$ be the closed convex hull in
$L^2(A,\tau)$ of the countable set
\[
 \{W_t(x)W_t(y)^*:t\in G\}.
\]
It has a unique vector $c(x,y)$ of least Hilbert norm.  To prove
measurability explicitly, enumerate its rational convex combinations
as $z_j(x,y)$ and put $d(x,y)=\inf_j\norm{z_j(x,y)}_2$.
These are measurable functions: bounded weakly measurable operator
fields are strongly measurable as vectors in the separable trace
Hilbert space.  For each $n$, choose the first $j$ with
$\norm{z_j(x,y)}_2^2<d(x,y)^2+2^{-n}$.  The choices are measurable.
For two such choices, the parallelogram identity and the lower bound
$d$ on the norm of their midpoint give
\[
 \norm{z_j-z_k}_2^2
 \le2(\norm{z_j}_2^2-d^2)+2(\norm{z_k}_2^2-d^2).
\]
They converge pointwise in $L^2$ to the measurable field $c$.

Every $z_j$ is an operator contraction with $\norm{z_j-I}\le r$.
These inequalities persist under bounded weak-operator limits, and
weak-operator compactness in the trace representation identifies any
such limit with the $L^2$ limit by its value on the trace vector.
Consequently $c(x,y)\in A$, $\norm{c(x,y)}\le1$ and
$\norm{c(x,y)-I}\le r$.  On bounded subsets, trace-Hilbert
measurability implies weak operator measurability in the trace
representation: test first on the dense vectors $p,q\in A\subset
L^2(A,\tau)$, where the coefficients are $\tau(q^*cp)$, and then
use the uniform operator bound.  Thus $c$ is also a measurable bounded
operator field.

The linear isometry
\[
 z\longmapsto W_g(x)zW_g(y)^*
\]
from the trace Hilbert space to itself sends
$C_{g^{-1}x,g^{-1}y}$ onto $C_{x,y}$.  Indeed the cocycle identity
sends its generator with index $t$ to the generator with index $gt$.
Uniqueness of the least-norm vector gives the asserted equivariance
for $c$.  Since $\norm{c-I}\le r<1$, its unitary polar part $a$ is
measurable, satisfies $\norm{a-I}\le2r$, and has the same
transformation law: polar decomposition respects multiplication on
either side by unitaries.  Finally
$\norm{W_g(x)W_g(y)^*-I}\le2\delta$ follows by comparing both
$W_g(x)$ and $W_g(y)$ with $u_g$.
\end{proof}

It remains to establish the rigidity of the diagonal action.
The bilateral-shift argument below proves the unitary-coefficient form
of the double-ergodicity theorem for Poisson boundaries
\cite{Kaimanovich2003}; see also
\cite[Theorem~2.1 and Remark~2.2, p.~7]{Bjorklund2014}.
Symmetry of the random-walk measure identifies the forward and
reflected boundary laws used in the proof.

\begin{lemma}[Double ergodicity with unitary coefficients]
\label{lem:double-ergodicity}
For the space $D$ constructed in \Cref{lem:boundary-mean}, let
$\sigma:G\to\U(K)$ be a unitary representation on a separable Hilbert
space.  If a bounded measurable map $f_0:D^2\to K$ satisfies
\[
 f_0(gx,gy)=\sigma(g)f_0(x,y)\quad(g\in G)
\]
almost everywhere, then $f_0$ is almost everywhere a constant vector
fixed by every $\sigma(g)$.
\end{lemma}

\begin{proof}
Use bilateral independent coordinates $(\omega_j)_{j\in\Z}$ with law
$\mu$, and let $\theta$ be the left shift, so
$(\theta\omega)_j=\omega_{j+1}$.  Define
\[
 b_+(\omega)=b(\omega_0,\omega_1,\ldots),\qquad
 b_-(\omega)=b(\omega_{-1}^{-1},\omega_{-2}^{-1},\ldots),\qquad
 f(\omega)=f_0(b_+(\omega),b_-(\omega)).
\]
Symmetry of $\mu$ makes $b_+$ and $b_-$ independent with common law
$\nu$.  The identities $b\circ S=b$ and equivariance give
\[
 b_+(\theta\omega)=\omega_0^{-1}b_+(\omega),\qquad
 b_-(\theta\omega)=\omega_0^{-1}b_-(\omega).
\]
For the first identity apply $b\circ S=b$ to the forward sequence;
for the second, apply it to
$(\omega_0^{-1},\omega_{-1}^{-1},\ldots)$ and then use equivariance.
All these identities hold almost everywhere, including after every
integer shift.  Consequently
\begin{equation}\label{eq:boundary-bilateral-transport}
 f(\theta^n\omega)
 =\sigma(\omega_{n-1}^{-1}\cdots\omega_0^{-1})f(\omega)
 \quad(n\ge1).
\end{equation}

Let $f_m$ be the Hilbert-valued conditional expectation of $f$ on
coordinates $[-m,m]$.  Then $\eps_m=\norm{f-f_m}_2\to0$ by density
of finite-coordinate functions in the product $L^2$ space.  Fix
$n>2m+1$ and put
\begin{align*}
 Q&=\sigma(\omega_m^{-1}\cdots\omega_0^{-1}),\qquad
 A=Qf_m,\\
 P&=\sigma(\omega_{n-m-1}^{-1}\cdots\omega_{m+1}^{-1}),\\
 R&=\sigma(\omega_{n-1}^{-1}\cdots\omega_{n-m}^{-1}),\qquad
 B=R^{-1}f_m\circ\theta^n.
\end{align*}
Thus $A$, $P$ and $B$ depend respectively on the disjoint coordinate
sets $[-m,m]$, $[m+1,n-m-1]$ and $[n-m,n+m]$; they are independent.
Equation \eqref{eq:boundary-bilateral-transport} and unitarity yield
\begin{equation}\label{eq:boundary-independent-blocks}
 \norm{PA-B}_2\le2\eps_m.
\end{equation}
Set $a_0=\E A$ and $b_0=\E B$.  Independence and orthogonality of
centered Hilbert-valued random variables give
\[
 \E\norm{PA-B}^2
 =\E\norm{A-a_0}^2+\E\norm{B-b_0}^2
   +\E\norm{Pa_0-b_0}^2.
\]
For example, condition first on $P$: the centered terms $P(A-a_0)$
and $B-b_0$ have zero mean, are independent, and their squared norms
are unchanged by $P$.  It follows that
\[
 \norm{A-a_0}_2\le2\eps_m,\qquad
 \norm{B-b_0}_2\le2\eps_m.
\]
Undoing the first transport shows
$\norm{f-Q^{-1}a_0}_2\le3\eps_m$.
This approximation depends only on coordinates $[0,m]$.
Undoing the other transport shows
$\norm{f\circ\theta^n-Rb_0}_2\le3\eps_m$.
After shifting back by $n$, the approximating function depends only on
coordinates $[-m,-1]$.  Hence $f$ belongs both to the closed $L^2$
subspace of functions of the nonnegative coordinates and to that of
functions of the negative coordinates.  The two sigma algebras are
independent, so their common subspace consists of constants: if $f$ is
measurable with respect to the second one, its conditional expectation
on the first is its constant expectation, and this conditional
expectation also equals $f$.  Write $f=v$ almost everywhere.
The joint law of $(b_+,b_-)$ is $\nu\times\nu$, so $f_0=v$ almost
everywhere.  Its equivariance and nonsingularity imply
$\sigma(g)v=v$ for every $g$.
\end{proof}
